% !PRINTED 4
\weberpage{4}
%<<CONT
Primzahl oder eine höhere Potenz einer der Primzahlen $a,b,c,\ldots$
in $m$ aufgehen müsste.

\begin{center}
\rule{0.18\linewidth}{0.4pt}
\end{center}

Wir geben nun einen Ueberblick über die in der Mathematik nothwendigen
und allmählich eingeführten Erweiterungen des Zahlenbegriffes.

Wir verstehen unter einer \gesperrt{Mannigfaltigkeit oder Menge}, oder
dem abkürzenden Zeichen $\mathfrak M$, ein System von Objecten oder
Elementen irgend welcher Art, das so in sich abgegrenzt und vollendet
ist, dass von jedem beliebigen Object vollkommen bestimmt ist, ob es
zu dem System gehört oder nicht, gleichviel, ob wir im Stande sind, in
jedem besonderen Falle die Entscheidung wirklich zu treffen oder nicht.

Eine Menge heisst \gesperrt{geordnet}, wenn von irgend zwei
unterschiedenen ihrer Elemente immer ein in sich vollkommen bestimmtes
als das grössere gilt, und zwar so, dass aus $a>b$, $b>c$ stets $a>c$
folgt. Ist $a>b$ und $b>c$ oder $a>b>c$, so sagen wir, dass $b$ zwischen
$a$ und $c$ liegt.% [source-original sic] Der Druck wiederholt den absteigenden Fall; die logisch zu erwartende aufsteigende Alternative wird nicht stillschweigend eingesetzt.

Die natürlichen Zahlen bilden eine geordnete Menge; zwischen zwei auf
einander folgenden ihrer Elemente liegt kein weiteres Element. Eine
solche Mannigfaltigkeit heisst eine \gesperrt{discrete}. Eine geordnete
Menge von der Eigenschaft, dass zwischen je zwei Elementen immer noch
andere Elemente gefunden werden, heisst \gesperrt{dicht}. Eine dichte
Menge kann man bilden, wenn man die natürlichen Zahlen in Paaren
zusammenfasst, und diese Paare als Elemente einer Menge auffasst. Diese
Paare sollen \gesperrt{Brüche} genannt und mit $m:n$ oder $m/n$
bezeichnet werden, und zwei solche Brüche $m:n$ und $m':n'$ werden
einander gleich gesetzt, wenn $mn'=nm'$ ist. Fasst man alle unter
einander gleichen Brüche zu einem Element zusammen, so erhält man eine
Mannigfaltigkeit, die geordnet ist, wenn man noch festsetzt, dass $m:n$
grösser als $m':n'$ ist, wenn $mn'>nm'$ ist. Dass diese Mannigfaltigkeit
dicht ist, sieht man so ein: sind $\mu=m:n$, $\mu'=m':n'$ zwei Brüche
und $\mu>\mu'$, so kann man, wenn $h$ eine willkürliche Zahl ist,
\[
 \mu={hmn'\over hnn'},
 \qquad
 \mu'={hm'n\over hnn'}
\]
%CONT>>
